\chapter{Реализованные шаблоны запросов}

В этом приложении фиксируются основные шаблоны, уже реализованные в репозитории. Они включены сюда, поскольку задают операциональную форму запросы, используемых на протяжении пилотного эксперимента.

\section{Запрос для генерации эталонных ответов}

Запрос, используемый для поиска эталонных ответов и генерацию кандидатов:
\begin{quote}
 \ttfamily
 Напиши шутка используя следующих ключевое слово(s): \{ключевые слова\}
\end{quote}

Его главное преимущество --- согласованность. Одна и та же текстовая форма используется для запроса для векторного представления во время поиск и для последующей генерацию кандидатов, что уменьшает несоответствие между построение набор данных и оценка. Главный недостаток состоит в том, что шаблон намеренно общий и поэтому не кодирует ни теорию юмора, ни информацию об аудитории.

\section{Инструкции для запросов для векторного поиска}

Для адаптации моделью векторных представлений к поиск подзадачи используются два коротких инструкция шаблоны:
\begin{quote}
 \ttfamily
 Instruct: Given ключевое слово для шутка, извлекать шутки would совпадение ключевое слово.\\
 Query: \{ключевое слово\}
\end{quote}

\begin{quote}
 \ttfamily
 Instruct: Given запрос для шутка, извлекать шутки что would ответ запрос.\\
 Query: \{запрос\}
\end{quote}

Эти шаблоны важны, потому что пилотный не использует исходные строки ключевых слов как векторных представлений напрямую. Он кодирует их внутри инструкции, ориентированные на поиск, что делает векторное пространство семантического намерения более специализированным под задачу.

\section{Оценочный запрос}

Системная инструкция для попарная оценка просит оценщик сравнить две шутки, написанные для одного запроса, и выбрать ровно одного победитель. Инструкция явно не поощряет вознаграждать длину или стилистическую выверенность, если они не улучшают юмор. Такой дизайн намеренно прост и должен рассматриваться как базовая модель протокол оценки, а не как финальная рубрика, согласованная с человеческими оценками.

\chapter{Дополнительные замечания о программной архитектуре}

Репозиторий организован как набор асинхронный процессы:
\begin{itemize}
 \item \texttt{JokesPipeline}: скачивает и предобрабатывает сырые корпуса;
 \item \texttt{EmbeddingsPipeline}: вычисляет плотные векторные представления шуток;
 \item \texttt{KeywordsPipeline}: извлекает значимые ключевые слова через сходство и MMR;
 \item \texttt{ReferencesPipeline}: строит наборы данных вида «запрос — эталонные ответы» через векторный поиск;
 \item \texttt{CandidatesPipeline}: генерирует шутки по запросам;
 \item \texttt{EvaluationPipeline}: выполняет попарное оценивание и агрегацию рейтинга.
\end{itemize}

Эта архитектура важна для финального диплома, потому что она отделяет построение данных от сравнения моделей. В результате последующий этап обучения MRVF можно добавить без переписывания остальной инфраструктуры. Отсутствующий компонент --- именно тот, который теперь становится основной целью будущей работы: процесс обучения, который потребляет построенный запрос-набор эталонных ответов и оптимизирует формальную целевую функцию, введенную в Chapter 2.
